\documentclass[11pt,a4paper]{article}
\usepackage[utf8]{inputenc}
\usepackage[vietnamese]{babel}
\usepackage{amsmath,amssymb}
\usepackage{geometry}
\geometry{a4paper, total={180mm,260mm}, left=15mm, right=15mm, top=18mm, bottom=18mm}

\usepackage{tikz}
\usetikzlibrary{shapes.geometric, arrows.meta, positioning, calc}
\tikzset{
  arr/.style={-{Stealth[scale=0.8]}},
  every node/.append style={align=center}
}
\usepackage{tcolorbox}
\tcbuselibrary{skins,breakable}
\usepackage{fontawesome}
\usepackage{fancyhdr}
\usepackage{enumitem}

\pagestyle{fancy}
\fancyhf{}
\renewcommand{\headrulewidth}{0.5pt}
\renewcommand{\footrulewidth}{0.5pt}
\fancyhead[L]{\small\bfseries\color{violet!85!black}\faBook\ SỔ TAY KIẾN THỨC TOÁN 2}
\fancyhead[R]{\small\bfseries\color{violet!85!black}Lớp Toán Thầy Lĩnh -- loptoanthaylinh.vn}
\fancyfoot[C]{\small\color{gray}-- Sổ tay kiến thức: Bài 15. Ki-lô-gam (kg) --}

\setlength{\parindent}{0pt}
\setlength{\parskip}{0pt}

\begin{document}

% --- BANNER BÀI HỌC ---
\begin{tcolorbox}[
  enhanced,
  colback=violet!8!white,
  colframe=violet!70!black,
  arc=3mm, boxrule=1.2pt,
  left=3mm, right=3mm, top=2mm, bottom=2mm,
  halign=center
]
  {\LARGE\bfseries\color{violet!90!black}\faBookmark\ BÀI 15: KI-LÔ-GAM (KG)}\\[1mm]
  {\itshape\footnotesize\color{black!75} Đơn vị đo khối lượng ki-lô-gam, dụng cụ cân và các phép tính kèm đơn vị trên 1 trang A4}
\end{tcolorbox}

\vspace{-2mm}

\noindent
\begin{minipage}[t]{0.49\textwidth}\vspace{0pt}
\begin{tcolorbox}[
  enhanced,
  colback=violet!4!white,
  colframe=violet!75!black,
  arc=2mm, boxrule=0.8pt,
  left=2mm, right=2mm, top=1.8mm, bottom=1.8mm,
  title={\footnotesize\bfseries\color{white}\faKey\ 1. Đơn vị đo ki-lô-gam}
]
{\footnotesize
\begin{itemize}\setlength{\itemsep}{1pt}\setlength{\leftskip}{-2.5mm}
  \item \textbf{Ki-lô-gam} là một đơn vị đo khối lượng phổ biến trong đời sống.
  \item Ki-lô-gam viết tắt là \textbf{kg}.
  \item \textbf{Dụng cụ đo:} Cân đĩa, cân đồng hồ, cân điện tử, cân sức khỏe,...
  \item \textbf{Cân thăng bằng (Cân đĩa):}\\
  Khi kim ở chính giữa (hai đĩa cân ngang bằng nhau), khối lượng vật ở đĩa bên này \textbf{bằng} khối lượng quả cân ở đĩa bên kia.
\end{itemize}
}
\end{tcolorbox}
\end{minipage}
\hfill
\begin{minipage}[t]{0.49\textwidth}\vspace{0pt}
\begin{tcolorbox}[
  enhanced,
  colback=teal!4!white,
  colframe=teal!75!black,
  arc=2mm, boxrule=0.8pt,
  left=2mm, right=2mm, top=1.8mm, bottom=1.8mm,
  title={\footnotesize\bfseries\color{white}\faSliders\ 2. Phép tính với ki-lô-gam}
]
{\footnotesize
\begin{itemize}\setlength{\itemsep}{1pt}\setlength{\leftskip}{-2.5mm}
  \item Khi cộng hoặc trừ các số đo khối lượng, ta thực hiện tính toán với các số như bình thường, rồi \textbf{viết thêm kí hiệu kg vào sau kết quả}.
  \item \textbf{Ví dụ:}
  \begin{itemize}\setlength{\itemsep}{0.5pt}\setlength{\leftskip}{-2.5mm}
    \item $15\text{ kg} + 8\text{ kg} = \mathbf{23\text{ kg}}$.
    \item $35\text{ kg} - 12\text{ kg} = \mathbf{23\text{ kg}}$.
    \item $20\text{ kg} + 10\text{ kg} - 5\text{ kg} = \mathbf{25\text{ kg}}$.
  \end{itemize}
  \item \textbf{So sánh:} Số nào lớn hơn thì khối lượng nặng hơn ($18\text{ kg} > 15\text{ kg}$).
\end{itemize}
}
\end{tcolorbox}
\end{minipage}

\vspace{-1.5mm}
\begin{tcolorbox}[
  enhanced,
  colback=orange!4!white,
  colframe=orange!80!black,
  arc=2mm, boxrule=0.8pt,
  left=2mm, right=2mm, top=1.5mm, bottom=1.5mm,
  title={\footnotesize\bfseries\color{white}\faSitemap\ 3. Sơ đồ tư duy: Bản chất đo khối lượng bằng cân đĩa}
]
\centering
\begin{tikzpicture}[x=1.2cm, y=0.8cm]
  % Cột cân
  \draw[thick, black!80] (3,0) -- (3,2);
  \draw[thick, fill=black!70] (2.8,0) rectangle (3.2,0.2);
  \draw[thick, fill=black!70] (2.9,1.9) circle (3pt);
  % Đòn cân thăng bằng
  \draw[very thick, violet!80!black] (1,1.8) -- (5,1.8);
  % Đĩa trái
  \draw[thick, teal!70!black] (1,1.8) -- (0.6,1) -- (1.4,1) -- cycle;
  \node[rectangle, draw=teal!80!black, fill=teal!20, font=\tiny\bfseries] at (1,1.2) {Túi gạo};
  % Đĩa phải
  \draw[thick, orange!70!black] (5,1.8) -- (4.6,1) -- (5.4,1) -- cycle;
  \node[rectangle, draw=orange!80!black, fill=orange!20, font=\tiny\bfseries] at (5,1.2) {Quả cân $2\text{ kg}$};
  % Kết luận
  \node[font=\scriptsize\bfseries, red!85!black] at (3, -0.4) {Kim chỉ ở giữa $\implies$ Khối lượng túi gạo $= 2\text{ kg}$};
\end{tikzpicture}
\end{tcolorbox}

\vspace{-1.5mm}
\noindent
\begin{minipage}[t]{0.49\textwidth}\vspace{0pt}
\begin{tcolorbox}[
  enhanced,
  colback=green!4!white,
  colframe=green!60!black,
  arc=2mm, boxrule=0.8pt,
  left=2mm, right=2mm, top=1.8mm, bottom=1.8mm,
  title={\footnotesize\bfseries\color{white}\faCheckCircle\ 4. Ví dụ thực hành chuẩn mực}
]
{\footnotesize
\textbf{Bài toán mẫu:}\\
Bao gạo to nặng $25\text{ kg}$, bao gạo bé nhẹ hơn bao gạo to $8\text{ kg}$. Hỏi bao gạo bé nặng bao nhiêu ki-lô-gam?\\
\textit{Bài giải:}\\
Bao gạo bé nặng số ki-lô-gam là:
$$25 - 8 = 17\text{ (kg)}$$
Đáp số: $17\text{ kg}$.
}
\end{tcolorbox}
\end{minipage}
\hfill
\begin{minipage}[t]{0.49\textwidth}\vspace{0pt}
\begin{tcolorbox}[
  enhanced,
  colback=red!4!white,
  colframe=red!75!black,
  arc=2mm, boxrule=0.8pt,
  left=2mm, right=2mm, top=1.8mm, bottom=1.8mm,
  title={\footnotesize\bfseries\color{white}\faExclamationTriangle\ 5. Mẹo tránh bẫy trắc nghiệm}
]
{\footnotesize
\begin{itemize}\setlength{\itemsep}{1pt}\setlength{\leftskip}{-2.5mm}
  \item \textbf{Bẫy quên ghi tên đơn vị "kg":}\\
  Kết quả phép tính bắt buộc phải có tên đơn vị: $15\text{ kg} + 5\text{ kg} = 20\text{ kg}$, ghi $20$ là bị trừ điểm.
  \item \textbf{Cân đĩa có quả cân ở cả 2 bên:}\\
  Nếu đĩa trái có túi đường và quả cân $1\text{ kg}$, đĩa phải có quả cân $5\text{ kg}$ thăng bằng:\\
  $\implies$ Túi đường nặng: $5 - 1 = 4\text{ kg}$.
\end{itemize}
}
\end{tcolorbox}
\end{minipage}


\end{document}
